\subsection{CENTRE OF THE ENVELOPING ALGEBRA}

In this number, we choose a basis B of R. Let \(R_{+}\) be the set of positive roots relative to B. Let \(\rho = \frac{1}{2} \sum_{\alpha \in R_{+}} \alpha\) , and \(\delta\) the automorphism of the algebra \(\mathbf{S}(\mathfrak{h})\) that takes every \(x \in h\) to \(x - \rho(x)\) , and hence the polynomial function p on \(h^{*}\) to the function \(\lambda \mapsto p(\lambda - \rho)\) .

THEOREM 2. Let U be the enveloping algebra of g, Z its centre, V ⊂ U the enveloping algebra of h (identified with S(h)), U⁰ the commutant of V in U, \(\varphi\) the Harish-Chandra homomorphism (§6, no. 4) from \(U^{0}\) to V relative to B. Let \(\mathbf{S}(\mathfrak{h})^{\mathrm{W}}\) be the set of elements of \(\mathbf{S}(\mathfrak{h})\) invariant under the action of W. Then \((\delta\circ\varphi)|Z\) is an isomorphism from Z to \(\mathbf{S}(\mathfrak{h})^{\mathrm{W}}\) , independent of the choice of B.

\begin{enumerate}
\def\labelenumi{\alph{enumi})}
\tightlist
\item
  Let \(P_{++}\) be the set of dominant weights of R, \(w \in W\) , \(\lambda \in P_{++}\) , \(\mu = w\lambda\) . Then \(Z(\mu - \rho)\) is isomorphic to a submodule of \(Z(\lambda - \rho)\) (§6, no. 3, Cor. 2 of Prop. 6), and \(\varphi(u)(\lambda - \rho) = \varphi(u)(\mu - \rho)\) for all \(u \in Z\) (§6, no. 4, Prop. 7). Thus, the polynomial functions \((\delta \circ \varphi)(u)\) and \((\delta \circ \varphi)(u) \circ w\) on \(h^{*}\) coincide on \(P_{++}\) . But \(P_{++}\) is dense in \(h^{*}\) in the Zariski topology: this can be seen by identifying \(h^{*}\) with \(k^{B}\) by means of the basis consisting of the fundamental weights \(\varpi_{\alpha}\) , and by applying Prop. 9 of Algebra, Chap. IV, §2, no. 3. Hence
\end{enumerate}

\[
(\delta \circ \varphi) (u) = (\delta \circ \varphi) (u) \circ w,
\]

which proves that \((\delta \circ \varphi)(\mathbf{Z}) \subset \mathbf{S}(\mathfrak{h})^{\mathrm{W}}\) .

\begin{enumerate}
\def\labelenumi{\alph{enumi})}
\setcounter{enumi}{1}
\tightlist
\item
  Let \(\eta\) be the isomorphism from \(\mathrm{I}(\mathfrak{g})\) to \(\mathbf{S}(\mathfrak{h})^{\mathrm{W}}\) defined in no. 3, Cor. 2 of Th. 1. Consider the canonical isomorphism from the g-module U to the g-module \(\mathbf{S}(\mathfrak{g})\) (Chap. I, §2, no. 8), and let \(\theta\) be its restriction to Z. Then \(\theta(\mathrm{Z}) = \mathrm{I}(\mathfrak{g})\) . Let z be an element of Z with filtration \(\leq f\) in U.
\end{enumerate}

\[
\begin{array}{c c c} \mathrm{Z} & \stackrel {{\theta}} {{\longrightarrow}} & \mathrm{I} (\mathfrak {g}) \\ \varphi \Big \downarrow & & \Big \downarrow^ {\eta} \\ \mathbf {S} (\mathfrak {h}) & \stackrel {{\delta}} {{\longrightarrow}} & \mathbf {S} (\mathfrak {h}). \end{array}
\]

Introduce the notations of §6, no. 4, and put

\[
z = \sum_ {\sum q _ {i} + \sum m _ {i} + \sum p _ {i} \leq f} \lambda_ {(q _ {i}), (m _ {i}), (p _ {i})} u ((q _ {i}), (m _ {i}), (p _ {i})).
\]

Let \(v((q_i),(m_i),(p_i))\) be the monomial \(X_{-\alpha_1}^{q_1}\ldots X_{-\alpha_n}^{q_n}H_1^{m_1}\ldots H_l^{m_l}X_{\alpha_1}^{p_1}\ldots X_{\alpha_n}^{p_n}\) calculated in \(\mathbf{S}(\mathfrak{g})\) . Denoting by \(\mathbf{S}_d(\mathfrak{g})\) the sum of the homogeneous components of \(\mathbf{S}(\mathfrak{g})\) of degrees \(0,1,\ldots,d\) , we have

\[
\theta (z) \equiv \sum_ {\sum q _ {i} + \sum m _ {i} + \sum p _ {i} = f} \lambda_ {(q _ {i}), (m _ {i}), (p _ {i})} v ((q _ {i}), (m _ {i}), (p _ {i})) \pmod {\mathbf {S} _ {f - 1} (\mathfrak {g})}
\]

SO

\[
(\eta \circ \theta) (z) \equiv \sum_ {\sum m _ {i} = f} \lambda_ {(0), (m _ {i}), (0)} v ((0), (m _ {i}), (0)) \pmod {\mathbf {S} _ {f - 1} (\mathfrak {h})}
\]

and consequently

\[
(\eta \circ \theta) (z) \equiv \varphi (z) \pmod {\mathbf {S} _ {f - 1} (\mathfrak {h})}.\tag{3}
\]

\begin{enumerate}
\def\labelenumi{\alph{enumi})}
\setcounter{enumi}{2}
\tightlist
\item
  We show that \(\delta \circ \varphi : \mathbf{Z} \to \mathbf{S}(\mathfrak{h})^{\mathrm{W}}\) is bijective. The canonical filtrations on U and \(\mathbf{S}(\mathfrak{g})\) induce filtrations on Z, \(\mathrm{I}(\mathfrak{g})\) and \(\mathbf{S}(\mathfrak{h})^{\mathrm{W}}\) , and \(\theta, \eta\) are compatible with these filtrations, so that \(\mathrm{gr}(\eta\circ\theta)\) is an isomorphism from the vector space \(\mathrm{gr}(\mathbf{Z})\) to the vector space \(\mathrm{gr}(\mathbf{S}(\mathfrak{h})^{\mathrm{W}})\) . By (3), \(\mathrm{gr}(\varphi)=\mathrm{gr}(\eta\circ\theta)\) , and it is clear that \(\mathrm{gr}(\delta)\) is the identity. Hence \(\mathrm{gr}(\delta\circ\varphi)\) is bijective, so
\end{enumerate}

\[
\delta \circ \varphi : Z \to \mathbf {S} (\mathfrak {h}) ^ {W}
\]

is bijective (Commutative Algebra, Chap. III, §2, no. 8, Cor. 1 and 2 of Th. 1). d) Recall the notations in a). Let E be a simple g-module of highest weight \(\lambda\) , and \(\chi\) its central character (§6, no. 1, Def. 2). Let \(\varphi'\) and \(\delta'\) be the homomorphisms analogous to \(\varphi\) and \(\delta\) relative to the basis \(w(\mathrm{B})\) . The highest weight of E relative to \(w(\mathrm{B})\) is \(w(\lambda)\) . By §6, no. 4, Prop. 7,

\[
\varphi (u) (\lambda) = \chi (u) = \varphi^ {\prime} (u) (w \lambda)
\]

for all \(u\in \mathbf{Z}\) , so, by a),

\[
\begin{array}{c} (\delta \circ \varphi) (u) (w \lambda + w \rho) = (\delta \circ \varphi) (u) (\lambda + \rho) = \varphi (u) (\lambda) = \varphi^ {\prime} (u) (w \lambda) \\ = (\delta^ {\prime} \circ \varphi^ {\prime}) (u) (w \lambda + w \rho). \end{array}
\]

Thus, the polynomial functions \((\delta \circ \varphi)(u)\) and \((\delta' \circ \varphi')(u)\) coincide on \(w(\mathrm{P}_{++}) + w\rho\) , and hence are equal.

COROLLARY 1. For all \(\lambda \in \mathfrak{h}^*\) , let \(\chi_{\lambda}\) be the homomorphism \(z \mapsto (\varphi(z))(\lambda)\) from \(Z\) to \(k\) .

\begin{enumerate}
\def\labelenumi{(\roman{enumi})}
\item
  If \(k\) is algebraically closed, every homomorphism from \(Z\) to \(k\) is of the form \(\chi_{\lambda}\) for some \(\lambda \in \mathfrak{h}^*\) .
\item
  Let \(\lambda, \mu \in \mathfrak{h}^*\) . Then \(\chi_{\lambda} = \chi_{\mu}\) if and only if \(\mu + \rho \in \mathrm{W}(\lambda + \rho)\) .
\end{enumerate}

If \(k\) is algebraically closed, every homomorphism from \(\mathbf{S}(\mathfrak{h})^{\mathrm{W}}\) to \(k\) extends to a homomorphism from \(\mathbf{S}(\mathfrak{h})\) to \(k\) (Commutative Algebra, Chap. V, §1, no. 9, Prop. 22, and §2, no. 1, Cor. 4 of Th. 1), and every homomorphism from \(\mathbf{S}(\mathfrak{h})\) to \(k\) is of the form \(f \mapsto f(\lambda)\) for some \(\lambda \in \mathfrak{h}^*\) (Chap. VII, App. I, Prop. 1). Hence, if \(\chi\) is a homomorphism from \(Z\) to \(k\) , there exists (Th. 2) a \(\mu \in \mathfrak{h}^*\) such that, for all \(z \in Z\) ,

\[
\chi (z) = ((\delta \circ \varphi) (z)) (\mu) = (\varphi (z)) (\mu - \rho)
\]

hence (i).

Let \(\lambda, \mu \in \mathfrak{h}^*\) and assume that \(\chi_{\lambda} = \chi_{\mu}\) . Then, for all \(z \in \mathbf{Z}\) ,

\[
((\delta \circ \varphi) (z)) (\lambda + \rho) = (\varphi (z)) (\lambda) = \chi_ {\lambda} (z) = \chi_ {\mu} (z) = ((\delta \circ \varphi) (z)) (\mu + \rho);
\]

in other words, the homomorphisms from \(\mathbf{S}(\mathfrak{h})\) to k defined by \(\lambda + \rho\) and \(\mu + \rho\) coincide on \(\mathbf{S}(\mathfrak{h})^{\mathrm{W}}\) ; thus, assertion (ii) follows from Commutative Algebra, Chap. V, §2, no. 2, Cor. of Th. 2.

COROLLARY 2. Let E, E' be finite dimensional simple g-modules, and \(\chi, \chi'\) their central characters. If \(\chi = \chi'\) , E and E' are isomorphic.

Let \(\lambda, \lambda'\) be the highest weights of E, E'. By §6, no. 4, Prop. 7, \(\chi_{\lambda} = \chi = \chi' = \chi_{\lambda'}\) , so there exists \(w \in \mathrm{W}\) such that \(\lambda' + \rho = w(\lambda + \rho)\) . Since \(\lambda + \rho\) and \(\lambda' + \rho\) belong to the chamber defined by B, we have \(w = 1\) . Thus, \(\lambda = \lambda'\) , hence the corollary.

PROPOSITION 7. For any class \(\gamma\) of finite dimensional simple g-modules, let \(U_{\gamma}\) be the isotypical component of type \(\gamma\) of the g-module U (for the adjoint representation of g on U). Let \(\gamma_{0}\) be the class of the trivial g-module of dimension 1. Let \([U,U]\) be the vector subspace of U generated by the brackets of pairs of elements of U.

\begin{enumerate}
\def\labelenumi{(\roman{enumi})}
\item
  U is the direct sum of the \(\mathrm{U}_{\gamma}\) .
\item
  \(U_{\gamma_{0}} = Z\) , and \(\sum_{\gamma \neq \gamma_{0}} U_{\gamma} = [U, U]\) .
\item
  Let \(u \mapsto u^{\natural}\) be the projection of U onto Z defined by the decomposition \(\mathrm{U} = \mathrm{Z} \oplus [\mathrm{U}, \mathrm{U}]\) . If \(u \in \mathrm{U}\) and \(v \in \mathrm{U}\) , we have \((uv)^{\natural} = (vu)^{\natural}\) . If \(u \in \mathrm{U}\) and \(z \in \mathrm{Z}\) , we have \((uz)^{\natural} = u^{\natural}z\) .
\item
  Let \(\varphi\) be the Harish-Chandra homomorphism. Let \(\lambda \in \mathrm{P}_{++}\) , and let E be a finite dimensional simple \(\mathfrak{g}\) -module of highest weight \(\lambda\) . For all \(u \in \mathrm{U}\) , we have
\end{enumerate}

\[
\frac {1}{\dim \mathrm{E}} \mathrm{Tr} (u _ {\mathrm{E}}) = (\varphi (u ^ {\natural})) (\lambda).
\]

The \(\mathfrak{g}\) -module U is a direct sum of finite dimensional submodules. This implies (i).

It is clear that \(U_{\gamma_{0}} = Z\) . Let \(U'\) be a vector subspace of U defining a subrepresentation of class \(\gamma\) of the adjoint representation. Then either \([g, U'] = U'\) or \([g, U'] = 0\) . Thus, if \(\gamma \neq \gamma_{0}\) then \([g, U'] = U'\) , so \(\sum_{\gamma \neq \gamma_{0}} U_{\gamma} \subset [U, U]\) . On the other hand, if \(u \in U\) and \(x_{1}, \ldots, x_{n} \in g\) , then

\[
\begin{array}{c} [ x _ {1} \ldots x _ {n}, u ] = (x _ {1} \ldots x _ {n} u - x _ {2} \ldots x _ {n} u x _ {1}) + (x _ {2} \ldots x _ {n} u x _ {1} - x _ {3} \ldots x _ {n} u x _ {1} x _ {2}) \\ + \dots + (x _ {n} u x _ {1} \ldots x _ {n - 1} - u x _ {1} \ldots x _ {n}) \in [ \mathfrak {g}, \mathrm{U} ]. \end{array}
\]

Hence \([\mathrm{U},\mathrm{U}]\subset \left[\mathfrak{g},\sum_{\gamma}\mathrm{U}_{\gamma}\right] = \left[\mathfrak{g},\sum_{\gamma \neq \gamma_0}\mathrm{U}_\gamma \right]\subset \sum_{\gamma \neq \gamma_0}\mathrm{U}_\gamma .\) This proves (ii). Under these conditions, (iii) follows from Chap. I, §6, no. 9, Lemma 5.

Finally, let E, \(\lambda\) be as in (iv). Then

\[
\begin{array}{r l} & {\mathrm{Tr} (u _ {\mathrm{E}}) = \mathrm{Tr} ((u ^ {\natural}) _ {\mathrm{E}}) \qquad \text {since} u - u ^ {\natural} \in [ \mathrm{U}, \mathrm{U} ]} \\ & {\qquad = \mathrm{Tr} (\varphi (u ^ {\natural}) (\lambda). 1) \qquad (\S 6, \text {no.} 4, \text {Prop.} 7)} \\ & {\qquad = (\dim \mathrm{E}). \varphi (u ^ {\natural}) (\lambda).} \end{array}
\]

